
%\documentclass{beamer}
%\usepackage{calc}
%
%%\usepackage{moreverb} % tab control
%\usepackage{fancyvrb} % tab


\documentclass[10pt]{beamer}
	\usetheme{Dresden}
%	\usetheme{Copenhagen}
%	\usetheme{Marburg}
%	\usetheme{CambridgeUS}
%	\usetheme{Warsaw}
%	\usetheme{AnnArbor}

%	\usecolortheme{orchid}
	
%	\usecolortheme{}
%	\useinnertheme{}
%	\useoutertheme{}
%	\setbeamertemplate{bibliography item}[text]
\usepackage{beamerthemesplit}
%\usecolortheme{default}
%\usecolortheme{structure}
%\usepackage{beamerthemesplit} // Activate for custom appearance

\useoutertheme{split}
\usecolortheme{whale}
\useinnertheme{rounded} 
\usecolortheme{orchid}
\setbeamertemplate{blocks}[rounded]




% ~~~~~~~~~~~~~~~~~~~~~~~~~~~~~~~~~~~~~~~ V HB lib
\usepackage{../hbcommon/hbdatastructures} 
\usepackage{../hbcommon/hbpresentation} 
% ~~~~~~~~~~~~~~~~~~~~~~~~~~~~~~~~~~~~~~~ A




% ~~~~~~~~~~~~~~~~~~~~~~~~~~~~~~~~~~~~~~~ V specific	
\usepackage[normalem]{ulem} %strikethroug

\usepackage{tikz-qtree} % tree
\usepackage{adjustbox} % tikzpicture alignment

%\usepackage{moreverb} % tab control
\usepackage{fancyvrb} % tab

\usepackage{epstopdf} % eps 2 pdf
\epstopdfsetup{update} % only regenerate pdf files when eps file is newer

%\usepackage[ruled,vlined,linesnumbered]{algorithm2e}
\usepackage[ruled,vlined]{algorithm2e}

\newcommand{\hDummy}{dummy}


% beamer blocks
% https://tex.stackexchange.com/questions/347269/different-colors-for-different-types-of-blocks-in-beamer
% https://tex.stackexchange.com/questions/28760/custom-beamer-blocks-for-pros-and-cons/

% https://tex.stackexchange.com/questions/1858/stable-vertical-alignment-of-columns-in-beamer
% https://tex.stackexchange.com/questions/127027/vertical-alignment-of-tikz-picture-in-a-table-cell
% https://tex.stackexchange.com/questions/172434/drawing-vertical-hierarchical-n-ary-tree-in-tikz
%  ~~~~~~~~~~~~~~~~~~~~~~~~~~~~~~~~~~~~~~~ A



\begin{document}




% ~~~~~~~~~~~~~~~~~~~~~~~~~~~~~~~~~~~~~~ V
\begin{frame}[t]\frametitle{Max-Heap Property-1}

	\vspace{-0.5cm}
	\begin{columns}[T] % >>>>
	\begin{column}{.5\linewidth} % ++++ V
	{\footnotesize 
		\begin{definition}
			An binary tree  \hCode{T} has  \hDefined{Max-Heap Property} 
			iff
			\hCode{i.parent $\ge$ i} for all nodes $i$.
			A binary tree with max-heap property is called \hDefined{max-heap}.
		\end{definition}
		
	} %fsize
	\end{column} % ++++ A
	\begin{column}{.5\linewidth} % ++++ V
		\begin{example}
		{\scriptsize
			\input{../tree/input/input-tree-CLRS6-1.tex}
		} % fsize
		{\tiny 

% Table generated by Excel2LaTeX from sheet 'Table'
\begin{tabular}
{
|@{~}p{.5cm}@{}
|@{~}p{.3cm}@{~}	|@{~}p{.3cm}@{~}
|@{~}p{.3cm}@{~}	|@{~}p{.3cm}@{~}	|@{~}p{.3cm}@{~}	|@{~}p{.3cm}@{~}	
|@{~}p{.3cm}@{~}	|@{~}p{.3cm}@{~}	|@{~}p{.3cm}@{~}	|@{~}p{.3cm}@{~}	|
}
\hline
index & 1     & 2     & 3     & 4     & 5     & 6     & 7     & 8     & 9     & 10      \\
\hline
value & 16    & 14    & 10    & 8     & 7     & 9     & 3     & 2     & 4     & 1       \\
\hline
\end{tabular}%

			\begin{tabular}
			{|
			@{~}l@{~}
			||@{~}r@{~}
			||@{~}r@{~} |@{~}r@{~}
			||@{~}r@{~}|@{~}r@{~}|@{~}r@{~}|@{~}r@{~}
			||@{~}r@{~}|@{~}r@{~}|@{~}r@{~}
			|}
			\hline
			index	&1	&2	&3	&4	&5	&6	&7	&8	&9	&10\\	
			\hline
			value	&16	&14	&10	&8	&7	&9	&3	&2	&4	&1\\
			\hline
			\end{tabular}

			\begin{tabular}
			{| l	|| r| r| r	| r	| r	| r	| r	| r	| r	| r	|}
			\hline
			index	&1	&2	&3	&4	&5	&6	&7	&8	&9	&10\\	
			\hline
			value	&16	&14	&10	&8	&7	&9	&3	&2	&4	&1\\
			\hline
			\end{tabular}

			\begin{tabular}
			{ | p{.2cm} ||p{.05cm}  | p{.1cm} |}
			index	&11	&22\\
			index	&33	&44\\
			\end{tabular}

			\begin{tabular}
%			{@{}p{.5cm}@{}p{.3cm}@{}p{.3cm}@{}}
			{|@{~}p{.5cm}@{}|@{~}p{.3cm}@{~}||@{~}p{.3cm}@{}|}
			index	&11	&22\\
			index	&33	&44\\
			\end{tabular}

			\begin{tabular}
%			{@{}p{.5cm}@{}p{.3cm}@{}p{.3cm}@{}}
			{
			|@{~}l@{~}
			||@{~}r@{~}
			||@{~}r@{~}|
			}
			index	&11	&22\\
			index	&33	&44\\
			\end{tabular}


		} % fsize

		\end{example}

	\end{column} % ++++ A
	\end{columns} % <<<<
\end{frame}
% ~~~~~~~~~~~~~~~~~~~~~~~~~~~~~~~~~~~~~~ A




% ~~~~~~~~~~~~~~~~~~~~~~~~~~~~~~~~~~~~~~ V
\begin{frame}[t]\frametitle{Tree Delete-2}

	\begin{tabular}
	{| l	|| r| r| r	| r	| r	| r	| r	| r	| r	| r	|}
	\hline
	index	&1	&2	&3	&4	&5	&6	&7	&8	&9	&10\\	
	\hline
	value	&16	&14	&10	&8	&7	&9	&3	&2	&4	&1\\
	\hline
	\end{tabular}

\setlength{\arrayrulewidth}{-1pt}
	\begin{tabular}
	{| l	|| r| r| r	| r	| r	| r	| r	| r	| r	| r	|}
	\hline
	index	&1	&2	&3	&4	&5	&6	&7	&8	&9	&10\\	
	\hline
	value	&16	&14	&10	&8	&7	&9	&3	&2	&4	&1\\
	\hline
	\end{tabular}
\setlength{\arrayrulewidth}{+1pt}

	\addtolength{\tabcolsep}{-1pt}
	\begin{tabular}
	{| l	|| r| r| r	| r	| r	| r	| r	| r	| r	| r	|}
	\hline
	index	&1	&2	&3	&4	&5	&6	&7	&8	&9	&10\\	
	\hline
	value	&16	&14	&10	&8	&7	&9	&3	&2	&4	&1\\
	\hline
	\end{tabular}
	\addtolength{\tabcolsep}{+1pt}

\end{frame}
% ~~~~~~~~~~~~~~~~~~~~~~~~~~~~~~~~~~~~~~ A




% ~~~~~~~~~~~~~~~~~~~~~~~~~~~~~~~~~~~~~~ V
\begin{frame}[t]\frametitle{Tree Delete-3}


%\newcolumntype{C}[1]{>{\centering\let\newline\\\arraybackslash\hspace{0pt}}m{#1}}
%%
%\begin{tabular}{C{1cm}C{1cm}}
%test for this & test for this
%\end{tabular}

	\begin{tabular}
	{| l	|| r| r| r	| r	| r	| r	| r	| r	| r	| r	|}
	\hline
	index	&1	&2	&3	&4	&5	&6	&7	&8	&9	&10\\	
	\hline
	value	&16	&14	&10	&8	&7	&9	&3	&2	&4	&1\\
	\hline
	\end{tabular}

	\addtolength{\tabcolsep}{-1pt}
	\begin{tabular}
	{| l	|| r| r| r	| r	| r	| r	| r	| r	| r	| r	|}
	\hline
	index	&1	&2	&3	&4	&5	&6	&7	&8	&9	&10\\	
	\hline
	value	&16	&14	&10	&8	&7	&9	&3	&2	&4	&1\\
	\hline
	\end{tabular}
	\addtolength{\tabcolsep}{+1pt}

	\addtolength{\tabcolsep}{-2pt}
	\begin{tabular}
	{| l	|| r| r| r	| r	| r	| r	| r	| r	| r	| r	|}
	\hline
	index	&1	&2	&3	&4	&5	&6	&7	&8	&9	&10\\	
	\hline
	value	&16	&14	&10	&8	&7	&9	&3	&2	&4	&1\\
	\hline
	\end{tabular}

	\addtolength{\tabcolsep}{0pt}
	\begin{tabular}
	{| l	|| r| r| r	| r	| r	| r	| r	| r	| r	| r	|}
	\hline
	index	&1	&2	&3	&4	&5	&6	&7	&8	&9	&10\\	
	\hline
	value	&16	&14	&10	&8	&7	&9	&3	&2	&4	&1\\
	\hline
	\end{tabular}

\addtolength{\tabcolsep}{+2pt}
	\begin{tabular}
	{| l	|| r| r| r	| r	| r	| r	| r	| r	| r	| r	|}
	\hline
	index	&1	&2	&3	&4	&5	&6	&7	&8	&9	&10\\	
	\hline
	value	&16	&14	&10	&8	&7	&9	&3	&2	&4	&1\\
	\hline
	\end{tabular}

\end{frame}
% ~~~~~~~~~~~~~~~~~~~~~~~~~~~~~~~~~~~~~~ A




% ~~~~~~~~~~~~~~~~~~~~~~~~~~~~~~~~~~~~~~ V
\begin{frame}[t]\frametitle{Tree Delete-4}

\begin{table}[t]
\setlength{\tabcolsep}{4pt}
\begin{tabular}
	{| l *{10}{| p{.2cm}} |}
	\hline
	index	&1	&2	&3	&4	&5	&6	&7	&8	&9	&10\\	
	\hline
	D	&16	&14	&10	&8	&7	&9	&3	&2	&4	&1\\
	\hline
\end{tabular}
\begin{tabular}
	{| l *{10}{| p{.2cm}} |}
	\hline
	index	&1	&2	&3	&4	&5	&6	&7	&8	&9	&10\\	
	\hline
	D	&16	&14	&10	&8	&7	&9	&3	&2	&4	&1\\
	\hline
\end{tabular}
\end{table}

\addtolength{\tabcolsep}{-2pt}
\begin{tabular}
	{| l *{10}{|r} |}
	\hline
	index	&1	&2	&3	&4	&5	&6	&7	&8	&9	&10\\	
	\hline
	valuA	&16	&14	&10	&8	&7	&9	&3	&2	&4	&1\\
	\hline
\end{tabular}
\addtolength{\tabcolsep}{+2pt}

\begin{tabular}
%	{| l *{10}{|p{.2cm}} |}
	{| l *{10}{|r} |}
	\hline
	index	&1	&2	&3	&4	&5	&6	&7	&8	&9	&10\\	
	\hline
	valuB	&16	&14	&10	&8	&7	&9	&3	&2	&4	&1\\
	\hline
\end{tabular}


\begin{tabular}
	{| l *{10}{|p{.2cm}} |}
	\hline
	index	&1	&2	&3	&4	&5	&6	&7	&8	&9	&10\\	
	\hline
	valuC	&16	&14	&10	&8	&7	&9	&3	&2	&4	&1\\
	\hline
\end{tabular}


%\begin{tabular}
%{|@{~} r*{6}{|c} @{} |}
%	\hline
%	BLA     & 1.0 & 1:6       & 78  & 468 & --  & 546  \\
%            & 3.0 & 1:2       & 183 & 366 & --  & 549  \\
%	\hline
%\end{tabular}

\end{frame}
% ~~~~~~~~~~~~~~~~~~~~~~~~~~~~~~~~~~~~~~ A




% ~~~~~~~~~~~~~~~~~~~~~~~~~~~~~~~~~~~~~~ V
\begin{frame}[t]\frametitle{aaa-5}
	\begin{columns}[T] % >>>>
	\begin{column}{.5\linewidth} % ++++ V
	bb
	{\tiny
	
\begin{table}[t]
\setlength{\tabcolsep}{4pt}
\begin{tabular}
	{| l *{10}{| p{.2cm}} |}
	\hline
	index	&1	&2	&3	&4	&5	&6	&7	&8	&9	&10\\	
	\hline
	D	&16	&14	&10	&8	&7	&9	&3	&2	&4	&1\\
	\hline
\end{tabular}
\begin{tabular}
	{| l *{10}{| p{.2cm}} |}
	\hline
	index	&1	&2	&3	&4	&5	&6	&7	&8	&9	&10\\	
	\hline
	E	&16	&14	&10	&8	&7	&9	&3	&2	&4	&1\\
	\hline
\end{tabular}
\end{table}

\begin{table}[t]
\addtolength{\tabcolsep}{-2pt}
\begin{tabular}
	{| l *{10}{|r} |}
	\hline
	index	&1	&2	&3	&4	&5	&6	&7	&8	&9	&10\\	
	\hline
	A	&16	&14	&10	&8	&7	&9	&3	&2	&4	&1\\
	\hline
\end{tabular}
\addtolength{\tabcolsep}{+2pt}
\end{table}

\begin{table}[t]
\begin{tabular}
%	{| l *{10}{|p{.2cm}} |}
	{| l *{10}{|r} |}
	\hline
	index	&1	&2	&3	&4	&5	&6	&7	&8	&9	&10\\	
	\hline
	B	&16	&14	&10	&8	&7	&9	&3	&2	&4	&1\\
	\hline
\end{tabular}
\end{table}

\begin{table}[t]
\setlength{\tabcolsep}{5pt}
\begin{tabular}
%	{| l *{10}{|p{.05cm}} |}
	{|p{.3cm} 
	!{\vrule width 1pt}
	*{10}{p{.05cm}|} 
	}
	\hline
	\textcolor{colorIndex}{index}	&\textcolor{colorIndex}{1}	&\textcolor{colorIndex}{2}	&\textcolor{colorIndex}{3}	&\textcolor{colorIndex}{4}	&\textcolor{colorIndex}{5}	&\textcolor{colorIndex}{6}	&\textcolor{colorIndex}{7}	&\textcolor{colorIndex}{8}	&\textcolor{colorIndex}{9}	&\textcolor{colorIndex}{10}\\
	\hline
%	index	&1	&2	&3	&4	&5	&6	&7	&8	&9	&10\\	
	\hIndex{index}	&\hIndex{1}	&\hIndex{2}	&\hIndex{3}	&\hIndex{4}	&\hIndex{5}	&\hIndex{6}	&\hIndex{7}	&\hIndex{8}	&\hIndex{9}	&\hIndex{10}\\
	\hline
	C	&16	&14	&10	&8	&7	&9	&3	&2	&4	&1\\
	\hline
\end{tabular}
\end{table}

\noindent
\begin{tabular}{|c!{\vrule width 4pt}c|c!{\vrule width 8pt}}
\hline
column1a & column2a & column 3a \\
%\ChangeRT{4pt}
column1b & column2b & column 3b \\
\hline
column1c & column2c & column 3c \\
%\ChangeRT{2pt}
column1d & column2d & column 3d \\
\hline
\end{tabular}



\( aaa\)
\( \mathrm{aaa}\)
	}  % fsize

	\end{column} % ++++ A
	\begin{column}{.5\linewidth} % ++++ V
	    	aaa
	\end{column} % ++++ A
	\end{columns} % <<<<
\end{frame}
% ~~~~~~~~~~~~~~~~~~~~~~~~~~~~~~~~~~~~~~ A




% ~~~~~~~~~~~~~~~~~~~~~~~~~~~~~~~~~~~~~~ V
\begin{frame}[t]\frametitle{aaa-6}
	\begin{columns}[T] % >>>>
	\begin{column}{.5\linewidth} % ++++ V
	    	aaa
		{\scriptsize 
		% Table generated by Excel2LaTeX from sheet 'Table'
\begin{table}[htbp]
    \begin{tabular}{lrrrrrrrrrr}
	\hline
\hIndex{index}	&\hIndex{1}	&\hIndex{2}	&\hIndex{3}	&\hIndex{4}	&\hIndex{5}	&\hIndex{6}	&\hIndex{7}	&\hIndex{8}	&\hIndex{9}	&\hIndex{10}\\
	\hline
value	&16	&14	&10	&8	&7	&9	&3	&2	&4	&1\\
	\hline
\textcolor{colorIndex}{index}	&\textcolor{colorIndex}{1}	&\textcolor{colorIndex}{2}	&\textcolor{colorIndex}{3}	&\textcolor{colorIndex}{4}	&\textcolor{colorIndex}{5}	&\textcolor{colorIndex}{6}	&\textcolor{colorIndex}{7}	&\textcolor{colorIndex}{8}	&\textcolor{colorIndex}{9}	&\textcolor{colorIndex}{10}\\
	\hline
value	&16	&14	&10	&8	&7	&9	&3	&2	&4	&1\\
    \end{tabular}%
\end{table}%
		} % fsize
	\end{column} % ++++ A
	\begin{column}{.5\linewidth} % ++++ V
	    	aaa
		\begin{itemize}
			\item
			aaa
		\end{itemize}
	\end{column} % ++++ A
	\end{columns} % <<<<
\end{frame}
% ~~~~~~~~~~~~~~~~~~~~~~~~~~~~~~~~~~~~~~ A




% ~~~~~~~~~~~~~~~~~~~~~~~~~~~~~~~~~~~~~~ V
\begin{frame}[t]\frametitle{aaa-7}
	\begin{columns}[T] % >>>>
	\begin{column}{.5\linewidth} % ++++ V
	    	aaa
		
\begin{tabular}{|l|l|l|l|}
\hline
 &  &  &  \\ 
\hline
 &  &  &  \\ 
\hline
 &  &  &  \\ 
\hline
 &  &  &  \\ 
\hline
\end{tabular}

\vspace{1ex}

{
\setlength\arrayrulewidth{2pt}
\begin{tabular}{|l|l|l|l|}
\hline
 &  &  &  \\ 
\hline
 &  &  &  \\ 
\hline
 &  &  &  \\ 
\hline
 &  &  &  \\ 
\hline
\end{tabular}
}

\begin{tabular}{|l|l|l|l|}
\hline
 &  &  &  \\ 
\hline
 &  &  &  \\ 
\hline
 &  &  &  \\ 
\hline
 &  &  &  \\ 
\hline
\end{tabular}

	\end{column} % ++++ A
	\begin{column}{.5\linewidth} % ++++ V
	    	aaa
		\begin{itemize}
			\item
			aaa
		\end{itemize}
	\end{column} % ++++ A
	\end{columns} % <<<<
\end{frame}
% ~~~~~~~~~~~~~~~~~~~~~~~~~~~~~~~~~~~~~~ A




% ~~~~~~~~~~~~~~~~~~~~~~~~~~~~~~~~~~~~~~ V
\begin{frame}[t]\frametitle{aaa-8}
		https://tex.stackexchange.com/questions/256731/changing-width-of-table-lines
	\begin{columns}[T] % >>>>
	\begin{column}{.5\linewidth} % ++++ V
		
	    	aaa
		{\tiny 
	
\begin{table}[t]
\setlength{\tabcolsep}{5pt}
\begin{tabular}
%	{| l *{10}{|p{.05cm}} |}
	{|p{.3cm} 
	!{\vrule width 1pt}
	*{10}{p{.05cm}|} 
	}
	\hline
	\textcolor{colorIndex}{index}	&\textcolor{colorIndex}{1}	&\textcolor{colorIndex}{2}	&\textcolor{colorIndex}{3}	&\textcolor{colorIndex}{4}	&\textcolor{colorIndex}{5}	&\textcolor{colorIndex}{6}	&\textcolor{colorIndex}{7}	&\textcolor{colorIndex}{8}	&\textcolor{colorIndex}{9}	&\textcolor{colorIndex}{10}\\
%	\hline
	\noalign{\hrule height 1pt}
	C	&16	&14	&10	&8	&7	&9	&3	&2	&4	&1\\
	\hline
\end{tabular}
\end{table}

		} % fsize


	\end{column} % ++++ A
	\begin{column}{.5\linewidth} % ++++ V
	    	aaa
		\begin{itemize}
			\item
			aaa
		\end{itemize}
	\end{column} % ++++ A
	\end{columns} % <<<<
\end{frame}
% ~~~~~~~~~~~~~~~~~~~~~~~~~~~~~~~~~~~~~~ A




%% ~~~~~~~~~~~~~~~~~~~~~~~~~~~~~~~~~~~~~~ V
%\begin{frame}[t]\frametitle{Tree Delete}
%	\begin{columns}[T] % >>>>
%	\begin{column}{.6\linewidth} % ++++ V
%		{\scriptsize
%			\input{input/input-alg-TreeDelete}
%		} % fsize	
%
%	\end{column} % ++++ A
%	\begin{column}{.4\linewidth} % ++++ V
%
%		{\small 
%		\input{input/input-tree-CLRS12-3}
%		} % fsize
%	
%		{\footnotesize 
%			\begin{tabular}{| l |}
%			\hline
%			Trace of {\scriptsize \Transplant(T, 13)}\\
%			\hline
%			~\\
%			\hline
%			\end{tabular}
%		} % fsize
%
%	\end{column} % ++++ A
%	\end{columns} % <<<<
%	
%\end{frame}
%% ~~~~~~~~~~~~~~~~~~~~~~~~~~~~~~~~~~~~~~ A



\end{document}